\documentclass[11pt]{article}
\usepackage[a4paper,margin=1.55cm]{geometry}
\usepackage{fontspec}
\usepackage{xeCJK}
\setmainfont{Noto Serif}
\setCJKmainfont[AutoFakeSlant=0.15]{Yu Gothic}
\setCJKsansfont{Yu Gothic}
\usepackage{amsmath,amssymb,mathtools}
\setlength{\parindent}{1em}
\setlength{\parskip}{0pt}
\setlength{\emergencystretch}{10em}
\allowdisplaybreaks
\raggedbottom
\newcommand{\D}{\Delta}
\newcommand{\ma}{\mathfrak a}
\newcommand{\mb}{\mathfrak b}
\newcommand{\mo}{\mathfrak o}
\newcommand{\mS}{\mathfrak S}
\newcommand{\mJ}{\mathfrak J}
\newcommand{\srcnumdisplay}[2]{\begin{equation*}\tag*{#1}\begin{aligned}#2\end{aligned}\end{equation*}}
\newcommand{\PXXIIIitem}[2]{\par\noindent\hangindent=2.2em\hangafter=1\makebox[2.2em][l]{#1.}#2\par}
\makeatletter
\renewcommand\footnotesize{\@setfontsize\footnotesize{8}{9.5}\selectfont}
\makeatother

\begin{document}

% P23-JA-U00 — sealed German lines 13507--13526
\section*{23. 代数的不変式と微分不変式}
\begin{center}
Jahresber. d. Deutschen Mathematiker-Vereinigung 32 (1923), S. 177--184
\end{center}

\begin{center}
{\Large\bfseries 代数的不変式と微分不変式。}\par
\vspace{0.45em}
1922年9月18日、ライプツィヒでの講演に基づく報告。\par
\vspace{1.0em}
ゲッティンゲンの \textsc{Emmy Noether} による。
\end{center}

\begingroup
\renewcommand{\thefootnote}{\arabic{footnote})}
\setcounter{footnote}{0}

代数的不変式および微分不変式の発展について報告するにあたり、私は両分野とも、\textsc{Hilbert} の言葉によれば素朴な時期と形式的な時期に続く批判的時期に話を限りたい。この批判的時期を特徴づける名は、代数的不変式については \textsc{Hilbert} 自身、微分不変式については \textsc{Riemann} である。内容の面からいえば、代数的不変式については、これらの問題をめぐって本質的に発達し、ここでその全き鋭さを発揮しえた代数学の算術的方法が、微分不変式については、形式的変分法の方法がそれを特徴づける。そこで私は、個々の対象を越える意義をもつ算術的方法について報告し、第二部では、まさに変分法の方法に即して達成される微分不変式の代数的不変式への従属だけを扱うことにする。

% P23-JA-U01 — sealed German lines 13527--13553
\textbf{1.} \emph{代数的不変式}の基礎には、周知のように、$x_1,\ldots,x_n$ における\emph{一般線型群}がある：
\srcnumdisplay{(1)}{x_i=\sum_k s_{ik}x'_k
\qquad\text{または略して}\qquad x=S(x').}
ここで $s_{ik}$ は不定元を表し、したがって $\D=|s_{ik}|$ は零でない。いま $f(a,x),g(b,x),\ldots$ を $x$ の形式で、それぞれ次数 $\alpha,\beta,\ldots$ をもち、不定元 $a,b,\ldots$ を係数とするものとする。このとき (1) によって $a,b,\ldots$ の\emph{誘導変換}が生じる。実際、
\[
\begin{aligned}
f(a,x)&=\sum a_{i_1\ldots i_n}x_1^{i_1}\cdots x_n^{i_n}
      =\sum{}' a'_{j_1\ldots j_n}{x'_1}^{j_1}\cdots {x'_n}^{j_n}
      =f(a',x'),\\
g(b,x)&=\sum b_{r_1\ldots r_n}x_1^{r_1}\cdots x_n^{r_n}
      =\sum{}' b'_{s_1\ldots s_n}{x'_1}^{s_1}\cdots {x'_n}^{s_n}
      =g(b',x')
\end{aligned}
\]
から、$a',b',\ldots$ は $a,b,\ldots$ の線型斉次関数となり、その係数は $s_{ik}$ に関してそれぞれ $\alpha,\beta,\ldots$ 次となる：
\srcnumdisplay{(2)}{a'=A_\alpha(a);\qquad b'=B_\beta(b)\ldots .}
ここで\emph{不変式}とは、変換 (1) および (2) に対する不変式をいう。すなわち、不定元 $a,b,\ldots,x,y,\ldots$（ここで $y=S(y')$ とも変換される）の整有理関数であって、変換を施すと置換行列式の冪を因子として自らを再現するものである：
\srcnumdisplay{(3)}{I(a',b',\ldots,x',y',\ldots)
=\D^\rho I(a,b,\ldots,x,y,\ldots).}
$I$ の係数は有理数、さらには有理整数であると仮定してよい。任意の数体 $P$ の係数をもつ不変式は、$P$ の係数を用いて、そのような特殊な不変式有限個によって線型的に表せるからである。同様に、$I$ が各系列について別々に斉次であると仮定しても制限にはならない。ここでも任意の不変式は、そのような特殊な不変式有限個の和として組み立てられるからである。

逆に、(1) と (2) は条件 (3) によって再び決定されることに注意しておく。\textsc{Ostrowski} と \textsc{Schur} が最近示したように、\footnote{A. \textsc{Ostrowski} と I. \textsc{Schur}, Über eine fundamentale Eigenschaft der Invarianten einer binären Form, Math. Zeitschr. 15 (1922), およびこれに続く \textsc{Ostrowski} の未刊行の研究。} 誘導変換は、各系列ごとに分離して働く線型変換の全体、すなわち、明示できる一定の例外を除いて、$x,y,\ldots$ を含まない任意の不変式 $I$ を許す変換の全体としても特徴づけられる。

二つの不変式の積は再び不変式であるが、和が不変式となるのは、重み、すなわち (3) の指数 $\rho$ が一致するとき、かつそのときに限られる。しかし行列式一の変換に制限すれば、任意の和も再び不変式となり、そのように制限された群に対するすべての不変式を尽くす。したがって整域、しかも\emph{斉次整域}が得られる。これは、ある多項式がその領域に属すれば、その斉次成分も一つずつ属するような領域であり、ここでは再び (3) に従う各系列ごとに斉次な不変式に帰着する。

ここで、不変式論、ひいては算術的代数学全体がその周囲で発展した根本問題が生じる：\emph{この整域は有限生成か、すなわち、任意の不変式を $I_1,\ldots,I_k$ によって整有理的に $I=G(I_1,\ldots,I_k)$ と表せるような有限の整基底 $I_1,\ldots,I_k$ をもつか。} この問題に対する \textsc{Hilbert} の解法は、整基底の問題をイデアル基底の問題へ還元することに基づく。二元の場合に対する \textsc{Gordan} の証明は、見通しのきかない記号計算のため拡張を許さず、\textsc{Mertens-Hilbert} の証明も二元形式を線型形式へ分解するため拡張を許さない。ここでは算術的な根本思想を強調しつつその論旨を手短に示し、次いで関連する三つの問題圏、\mbox{すなわち有限個}の手順による基底の実際の構成、部分群における有限性問題、および整数係数性を考慮した有限性問題へ進みたい。

% P23-JA-U02 — sealed German lines 13554--13574
\textbf{2.} （有限次）代数体における\emph{イデアルの定義}を思い起こそう。体の数からなる系 $\ma$ は、次を満たすときイデアルと呼ばれる。
\PXXIIIitem{1}{$\ma$ が、その体のすべての整数からなる領域 $\mo$ に属すること、}
\PXXIIIitem{2}{$\alpha$ と $\beta$ が属すれば、差 $\alpha-\beta$ も $\ma$ に属すること、}
\PXXIIIitem{3}{$\alpha$ が属すれば $\lambda\alpha$ も $\ma$ に属すること。ただし $\lambda$ は $\mo$ の任意の元とする。}
そして周知の定理によれば、すべてのイデアルは\emph{イデアル基底}をもち、$\ma=(\alpha_1,\ldots,\alpha_k)$ と書ける。すなわち、$\ma$ の元 $\alpha_1,\ldots,\alpha_k$ が有限個存在し、2. と 3. によって $\ma$ をそれらから導けるので、$\alpha=\lambda_1\alpha_1+\cdots+\lambda_k\alpha_k$ が $\ma$ の全要素を尽くす。\footnote{$k$ を二まで減らせることは、以下では問題にならない。}

このイデアルの定義は、$\mo$ が整域をなすということだけに基づく。したがってイデアルの概念もイデアル基底の概念も、$\mo$ を任意の整域に置き換えても文字どおり保たれる。

さて \textsc{Hilbert} の有限性証明は、次の三定理に分かれる。
\PXXIIIitem{1}{有理性領域を係数とする $n$ 個の不定元の全多項式の領域では、すべてのイデアルについてイデアル基底定理が成り立ち、したがって任意の系 $\mS$ についても成り立つ。}
\PXXIIIitem{2}{多項式からなる斉次整域 $\mJ$ が有限生成であるための必要十分条件は、$\mJ$ においてイデアル基底定理が成り立つことである。}
\PXXIIIitem{3}{不変式領域 $\mJ$ では、イデアル基底定理は次の強化形で成り立つ。全多項式の領域に関して作った任意のイデアル基底が同時に $\mJ$ のイデアル基底でもあり、したがって}
\[
 I=A_1I_1+\cdots+A_kI_k
 \qquad\text{とともに常に}\qquad
 I=i_1I_1+\cdots+i_kI_k
\]
が成り立つ。ここで $A$ は $a,b,\ldots,x,y,\ldots$ の多項式を、$i$ は不変式を表す。

1. の証明は、原理上、代数体の場合と同じ推論に基づく。どちらの場合も、イデアル基底の存在は、線型形式からなる加群についての一般定理の直接の帰結である。\footnote{これについては、代数関数の算術理論に関する私の比較報告、Jahresber. 28 (1920), S. 187 および Anm. 2), S. 188 を参照。} 1. から、多項式からなる任意の有限生成整域についてイデアル基底の存在が直ちに従い、斉次領域について逆も成り立つことは容易に分かる。特殊な不変式の性質が使われるのは 3. だけである。\textsc{Hilbert} では 3. は $\Omega$ 過程に関する周知の定理から従う。一方、近年 E. \textsc{Fischer} は、別の一般定理を基礎にすれば、一般線型群が各変換とともにその反傾変換、またはその共役複素変換を含むという性質だけから 3. を導けることを示した。\footnote{E. \textsc{Fischer}, Über die Endlichkeit der Invarianten. Göttinger Nachrichten 1915, およびライプツィヒ講演。}

% P23-JA-U03 — sealed German lines 13575--13583
\textbf{3.} \emph{基底の実際の構成}は、\emph{関数体}の理論を援用して有限性概念をさらに算術的に変形することに基づく。次が成り立つ。
\PXXIIIitem{4}{多項式からなる整域 $\mJ$ が有限生成であるための必要十分条件は、$\mJ$ の中に有限生成部分整域 $\mJ'$ があり、$\mJ$ が $\mJ'$ 上整であることである。このとき $\mJ'$ の任意の整基底に、これら整な量の基本系を加えれば $\mJ$ の基底となる。\footnote{\textsc{Hilbert} は（Über die vollen Invariantensysteme, Math. Annalen 42 (1893), § 1 および 2）、仮定した有限生成性からこれらの事実が従うことを示した。逆に、整依存性から有限生成性が従うことは、有理基底の存在から分かる（E. Noether, Körper und Systeme rationaler Funktionen, Math. Annalen 76 (1915), § 12）。}}
とくに不変式の場合のように、イデアル基底定理が強化形 3. で成り立つ斉次領域 $\mJ$ を考えるなら、そのような領域 $\mJ'$ は、常に存在する何らかの有限個の多項式 $I_1,\ldots,I_s$ を整基底として導かれるように選べる。しかも、それらが消えれば $\mJ$ の全多項式も消える。この事実は再びイデアル論の一般定理に基づく。
\PXXIIIitem{5}{多項式からなる任意のイデアル $\ma$ に対し、通常の有理整数 $r$ が存在して、$\ma$ のすべての零点で消える任意のイデアル $\mb$ について $\mb^r$ は $\ma$ で割り切れる。}

不変式の場合には、$I_1,\ldots,I_s$ の\emph{次数の上界}、したがってその個数 $s$ の上界も、基礎に置いた基本形式の次数だけに依存する形で決定できる。ここでも不変式の特殊な性質が使われる。すなわち、数値的に特殊化した基本形式の系が零でない不変式をもつための必要十分条件は、$\D=|s_{ik}|$ が変換後の係数上整であることである。

この上界から K. \textsc{Hentzelt}\footnote{私はこの研究を近く遺稿から刊行する予定である。} は、\emph{整基底に属する不変式の次数の上界}を求めた。この上界も基本形式の次数だけに依存する。すなわち彼は、5. の\emph{指数} $r$ について、$\ma$ のイデアル基底に属する多項式の個数と最高次数だけに依存し、それらの係数には依存せず、しかも与えられた数から有限個の手順で計算できる\emph{上界}を与えたのである。原理的には、これによって提起した問題は完全に解決された。ただし、その上界はきわめて高い。例えば三変数二次形式では、判別式、すなわち三次の不変式だけですでに $x$ を含まない全不変式の基底となるのに、\textsc{Hilbert} と \textsc{Hentzelt} の評価では数百万にまで達する。

% P23-JA-U04 — sealed German lines 13584--13587
\textbf{4.} 一般線型群の\emph{部分群}の不変式に関する有限性問題では、まだ部分的成果に達したにすぎない。\mbox{その方法}はほとんどもっぱら、\textsc{Study} が直交群について採った手続きに従い、部分群の不変式を一般線型群の同時不変式へ還元することに基づく。\textsc{Weitzenböck} は運動不変式とアフィン不変式について、\textsc{Deruyts} は半不変式と剪断不変式についてこれを実行した。\footnote{文献については、\textsc{Weitzenböck} の不変式論に関する百科事典項目 III E 1 を参照。} この方法がどこまで及ぶかはまだ未決である。2. の末尾で述べた \textsc{Fischer} の指摘も、一連の部分群について有限性問題を解く。とくに直交群については、これにより最も簡単な証明が得られる。しかし半不変式の例は、\textsc{Fischer} が近年指摘したように、\footnote{ライプツィヒ講演および Math. Zeitschrift。} 部分群では、有限生成性にもかかわらず、イデアル基底に関する強化形 3. が成り立つ必要はないことを示す。

この有限性問題を体論の意味で捉えると、\textsc{Hilbert} の\emph{相対的に整な関数}の問題、すなわち、有理関数体の多項式全体からなる整域が常に有限生成であるかという問題へ至る。この方向では、有限群の不変式に対し、際立った整基底を直接与えられるという、初等的に証明できる事実がある。それはガロア分解式の係数系である。\footnote{E. Noether, Der Endlichkeitssatz der Invarianten endlicher Gruppen. Math. Ann. 77 (1916).}

% P23-JA-U05 — sealed German lines 13588--13589
\textbf{5.} 不変式の有限性定理を\emph{整数係数性}に関して強化できるかという問題も、一般にはまだ解決されていない。確かに \textsc{Hilbert} が示したように、イデアル基底定理は全整数係数多項式の領域についても成り立ち、イデアル基底と整基底との関係 2. もそのまま保たれる。\footnote{この関係から、とくに、私が Math. Ann. 83 (1921) で展開した一般イデアル論の分解定理が、すべての有限生成整域について成り立つことが従う。} しかし、有限性には必要でない強化形 3. の証明は通用しない。\emph{二元不変式}の場合には、\emph{整数係数での有限生成性}を証明できる。\footnote{E. Noether, Die Endlichkeit des Systems der ganzzahligen Invarianten binärer Formen. Göttinger Nachrichten 1919.} その証明は、二元の場合の \textsc{Mertens-Hilbert} の有限性証明を整数係数性に関して強化したものだが、整数係数イデアル基底の定理だけでなく、二元形式を線型形式へ分解することも用いるため、より多くの不定元へ移すことはできない。これに対し、同様の基礎の上で、先に述べた証明を強化すれば、\emph{整数係数での有限生成性}が\emph{有限群の不変式}について得られる。ただしここでも、基底を実際に与えるのではなく存在証明だけが得られる。\footnote{11) で引用した論文の § 4。} ここでも、関数体論と一般イデアル論のさらなる発展だけがおそらく目標へ導くであろう。

% P23-JA-U06 — sealed German lines 13590--13622
\textbf{6.} ここで\emph{微分不変式}へ移り、序論で提起した、それを代数的不変式論へ還元する問題に進む。基礎となる群は、周知のように $x_i=x_i(y)$ によって次のように定義される：
\srcnumdisplay{(4)}{d x_i=\sum_k \frac{\partial x_i}{\partial y_k}d y_k;\qquad
 d\delta x_i=\sum_k \frac{\partial x_i}{\partial y_k}d\delta y_k+
\sum_{k,l}\frac{\partial^2 x_i}{\partial y_k\partial y_l}d y_k\,\delta y_l;\ \ldots .}
(4) によって誘導される変換へ移るには、任意の関数 $f(x,d x)=g(y,d y)$ を基礎に置けばよい。(4) に対する $\delta f,\delta^2f,\ldots$ の不変性を要求すると、$x$ と $d x$ による $f$ の導関数についての変換が得られる。例えば $d x$ に関して斉次な形式に限れば、
\[
 f(x,d x)=\sum a_{i_1\ldots i_n}(x)d x_1^{i_1}\cdots d x_n^{i_n}
 =\sum{}' a'_{j_1\ldots j_n}(y)d y_1^{j_1}\cdots d y_n^{j_n}
 =g(y,d y),
\]
$a'$ は $a$ に関して線型となり、$\frac{\partial a'}{\partial y}$ は $a$ と $\frac{\partial a}{\partial x}$ に関して線型となる、等々：
\srcnumdisplay{(5)}{a'=A_0(a),\qquad
\frac{\partial a'}{\partial y}=A_1\!\left(a,\frac{\partial a}{\partial x}\right),\qquad
\frac{\partial^2 a'}{\partial y^2}=A_2\!\left(a,\frac{\partial a}{\partial x},\frac{\partial^2 a}{\partial x^2}\right),\ldots}
そして問題は、変換 (4) と (5) に対する不変性である。このとき現れるすべての量
\[
a,\ \frac{\partial a}{\partial x},\ldots,\ d x,\ \delta x,\ d\delta x,\ldots;\ 
\frac{\partial x}{\partial y},\ \frac{\partial^2 x}{\partial y^2},\ldots
\]
は不定元とみなすべきであるから、各変換の行列式は確かに零でない。形式的に完成した理論を特殊な関数へ適用するときに初めて、ある領域で変換が一意に可逆なままであり、十分な数の微分商が存在するような関数へ制限する問題が生じる。これは、代数の場合、数値的特殊化に際して行列式 $|s_{ik}|$ が零でないと仮定しなければならないのとまったく同じである。

このように、現れるすべての引数を不定元とみなすとき、通常の代数的方法では直接扱えない特殊な線型群を考えることになる。この群では $d x$ の変換と $a'=A_0(a)$ の変換が (1) およびそこから誘導された (2) と一致する。そこで、一般に\emph{新しい引数を導入する}ことによって、変換 (4) と (5) を一般線型群 (1) と、それによって\emph{誘導される変換} (2) へ\emph{還元できる}か、という問題が生じる。

これは実際に可能である。高次微分 $d^2x,\delta d x,\ldots$ は、$f(x,d x)$ に属する変分問題から生じるラグランジ\nolinebreak[4]ュの変分導関数と、さらにそれらを極化して得られる結合（\textsc{Weyl} と \textsc{Schouten} のいう接続）で置き換えられる。これらはすべて $d x$ と同じ変換則に従い、その変換は一般線型群 (1) によって与えられる。一方、$\frac{\partial a}{\partial x},\frac{\partial^2a}{\partial x^2},\ldots$（より一般には、斉次と仮定した $f$ の $x$ と $d x$ による導関数）は、「高次変分の正規形」から、上記の $d^2x,\delta d x,\ldots$ の消去によって生じる一定の形式の係数と、それらの「共変微分」
\[
 [\Omega_i],\quad [\Omega_i^{(1)}],\quad [\Omega_i^{(2)}],\ldots\qquad (i=1,2,\ldots);
\]
で置き換えられる。そして $[\Omega_i^{(k)}]$ は (4) と (5) に対する不変式で $d x$ と $\delta x$ だけに依存するから、その係数は代数的に誘導される変換 (2) を満たす。\footnote{E. Noether, Invarianten beliebiger Differentialausdrücke. Göttinger Nachrichten. 1918.} この\emph{還元定理}により、微分不変式論は代数的不変式論に従属する。

$f(x,d x)$ が $d x$ に関して $\alpha$ 次の斉次形式であるとき、$[\Omega_i]$ の列は $[\Omega_\alpha]$ で終わる。二次形式では $[\Omega_2]$ は \textsc{Riemann} の曲率形式と一致し、ここで示した構成過程は、\textsc{Riemann} がラテン語の懸賞論文で与え、またそれとは独立に、\textsc{Riemann} の教授資格講演を受けて \textsc{Lipschitz} が展開した過程とまったく同じである。\footnote{7) で挙げた \textsc{Weitzenböck} の百科事典項目の第二部第28項にある、\textsc{Riemann} と \textsc{Lipschitz} に関する私の叙述も参照。} \textsc{Christoffel} と \textsc{Ricci} が二次微分形式について計算で示した還元定理の証明は、「\textsc{Riemann} 正規座標」の導入に基づく。この座標は一点から出る極値曲線を直線に変換する。これによって $x$ の任意の変換に正規座標の線型変換が対応するので、正規座標が一般線型群への移行をもたらす。

\textsc{Weyl} と \textsc{Schouten} に従い、群の定義に変分問題ではなく\emph{接続を基礎に置く}場合にも、このような還元定理があるかという問題が生じる。すなわち、ラグランジュの変分導関数に対応して、高次微分を消去する式 $\psi_i(d\delta x,d x,\delta x)$ を立て、二次微分形式との類比により、それ自体は不要な制限、すなわち $\psi_i$ が $d x$ に関して線型であるという制限を加える場合である（任意の斉次 $f(x,d x)$ では、極化したラグランジュの変分導関数は $d x$ に関して一次斉次にすぎない）。$\psi$ が $d x$ と同じ変換則に従うことを要求すると、その係数の変換が誘導され、それによって (5) と同様に群が定まる。最低階の不変式については、\textsc{Weitzenböck} が \textsc{Weyl} の場合に還元定理を計算で確認したが、一般的な構成法はまだない。\footnote{次を参照：\textsc{Weyl}, Raum, Zeit, Materie; \textsc{Schouten}, Math. Zeitschr. 13 (1922); \textsc{Weitzenböck}, Sitzungsber. d. Wiener Akademie, 1920 u. 1921.}

% P23-JA-U07 — sealed German lines 13623--13630
\begin{center}
（1922年10月26日受理。）
\end{center}
\endgroup

\clearpage
\setcounter{footnote}{0}

\end{document}
